% !TeX root = ../../Book3.tex
\section{Dirichlet 级数}
\label{b3:sec:2504}

一般幂级数以圆盘为收敛区域，Dirichlet 级数以右半平面为基本区域。原因是 \(n^{-s}=e^{-s\log n}\)，增大实部会同时压低全部高指标项。

\subsection{收敛横坐标}
\label{b3:subsec:250401}

\begin{definition}\label{b3:def:25-dirichlet-series}
形如 \(D_a(s)=\sum_{n\ge1}a(n)n^{-s}\) 的级数称为\term[dirichlet-jishu]{Dirichlet 级数}[Dirichlet series]。
\end{definition}
若级数在 \(s_0\) 收敛，令 \(b_n=a(n)n^{-s_0}\)。其尾部分和 \(\sum_{M<n\le N}b_n\) 对 \(N\ge M\) 一致趋零。对尾部应用 Abel 求和可知，对任意 \(\operatorname{Re}w>0\)，\(\sum b_nn^{-w}\) 收敛，并在 \(\operatorname{Re}w\ge\delta>0\)、\(|w|\le R\) 一致收敛。因此存在扩展实数 \(\sigma_c\)，使级数在 \(\operatorname{Re}s>\sigma_c\) 收敛且全纯，在 \(\operatorname{Re}s<\sigma_c\) 发散；边界可能随虚部而异。数 \(\sigma_c\) 称为\term[shoulian-hengzuobiao]{收敛横坐标}[abscissa of convergence]。

\subsection{绝对收敛}
\label{b3:subsec:250402}

同样定义\term[juedui-shoulian-hengzuobiao]{绝对收敛横坐标}[abscissa of absolute convergence] \(\sigma_a\)。总有 \(\sigma_c\le\sigma_a\le\sigma_c+1\)，按扩展实数理解。后一个不等式由在 \(s_0\) 收敛推出 \(a(n)n^{-s_0}\) 有界，再与 \(\sum n^{-1-\varepsilon}\) 比较得到。绝对收敛半平面的每个较小闭半平面上有统一主控，因此可以安全乘法和对素数重新分组。

\ProseParagraphBreak

对普通 \(\sum n^{-s}\)，两横坐标都为一；对交错级数，两者分别为零与一。横坐标为零的必要性可由负实部处单项不趋零看出。一个例子同时说明两横坐标可以不同，也说明一般差距上界一能达到。

\subsection{乘法}
\label{b3:subsec:250403}

\begin{proposition}\label{b3:prop:25-product-series}
若 \(D_a,D_b\) 在 \(\operatorname{Re}s=\sigma\) 都绝对收敛，则
\(D_a(s)D_b(s)=D_{a*b}(s)\) 在该直线上成立，且右侧绝对收敛。
\end{proposition}
\begin{proof}
双重绝对和为 \((\sum|a(m)|m^{-\sigma})(\sum|b(n)|n^{-\sigma})<\infty\)。因此可以按积 \(mn=r\) 分组，系数恰为有限和 \(\sum_{m\mid r}a(m)b(r/m)\)。同一个绝对和控制了右侧级数。
\end{proof}
仅仅知道两个级数逐点条件收敛，不能引用这个结论；重新分组的合法性需要另证。

\ProseParagraphBreak

在绝对收敛半平面内部还可逐项求任意次导数。设 \(\sigma_0>\sigma_a\)，再取 \(\delta>0\)。对每个固定整数 \(k\ge0\)，函数 \((\log t)^kt^{-\delta}\) 在 \(t\ge1\) 上有界。因此在 \(\operatorname{Re}s\ge\sigma_0+\delta\) 上，
\[
 \sum_n|a(n)|(\log n)^k n^{-\operatorname{Re}s}
 \le C_{k,\delta}\sum_n|a(n)|n^{-\sigma_0}<\infty.
\]
由一致收敛的导数定理，
\[
 D_a^{(k)}(s)=(-1)^k\sum_{n\ge1}a(n)(\log n)^k n^{-s}.
\]
增加对数权重会影响边界收敛，却不改变右半平面内部的全纯性。这个论证说明为什么下一节可以对 Euler 对数求导，也说明每次求导都应在离边界正距离的紧集上进行估计。

\ProseParagraphBreak

Dirichlet 级数也可视为离散测度的 Laplace 变换：将质量 \(a(n)\) 放在 \(\log n\)，则 \(n^{-s}=e^{-s\log n}\)。相应累积质量为 \(A(e^t)\)。在端点项消失的范围，变量代换 \(x=e^t\) 将部分求和写成
\[
 D_a(s)=s\int_0^\infty A(e^t)e^{-st}\,\mathrm dt.
\]
这项表示把整数上的乘法尺度变为实轴上的加法尺度。若已经知道 \(A(x)\) 的增长，积分容易估计；反过来，只知道变换在某点有极点，通常还不能恢复 \(A(x)\) 的渐近。第二十七章将增加非负性与整条边界的延拓条件，完成那个反向问题。

\subsection{系数卷积}
\label{b3:subsec:250404}

于是 \(D_d=\zeta^2\)、\(D_{\sigma_\alpha}(s)=\zeta(s)\zeta(s-\alpha)\)，最初分别在 \(\operatorname{Re}s>1\) 与 \(\operatorname{Re}s>\max(1,1+\operatorname{Re}\alpha)\) 由绝对收敛证明。\(\mathbf1*\varphi=\operatorname{id}\) 还给 \(D_\varphi(s)=\zeta(s-1)/\zeta(s)\)，但要作除法须先在此处证明分母非零，下一节的 Euler 乘积将提供这一事实。
